\chapter*{RESUME}
\addcontentsline{toc}{chapter}{RESUME}

Le diabète sucré constitue l'un des enjeux de santé publique les plus critiques du
21\textsuperscript{e} siècle, touchant plus de 500 millions de personnes à travers le
monde. Sa détection précoce est essentielle pour prévenir des complications graves et
réduire la charge des systèmes de santé. Dans ce contexte, les méthodes d'apprentissage
automatique supervisé offrent des perspectives prometteuses pour l'exploitation des données
cliniques à des fins diagnostiques.

Ce travail présente une évaluation comparative rigoureuse de trois algorithmes de machine
learning — les \textbf{Machines à Vecteurs de Support (SVM)}, les \textbf{Forêts Aléatoires
(Random Forests)} et \textbf{XGBoost} — appliqués à la classification du risque diabétique
sur la base de données \textit{Pima Indians Diabetes Database} (PIDD). Ce jeu de données,
composé de 768 patientes issues de la communauté Pima présentant une prévalence élevée de
diabète de type 2, regroupe huit variables cliniques clés telles que la glycémie plasmatique,
l'indice de masse corporelle, l'insulinémie, la pression artérielle et la fonction de
prédisposition héréditaire.

Dans un premier temps, les fondements théoriques de chacun des trois algoritmes sont
présentés : le formalisme de la maximisation de la marge et l'astuce du noyau pour les SVM,
les mécanismes de bagging et de sélection aléatoire des caractéristiques pour les forêts
aléatoires, ainsi que le gradient boosting régularisé et l'approximation de Taylor de
second ordre propres à XGBoost. Un état de l'art de l'application de ces méthodes à la
détection du diabète est ensuite conduit, soulignant les limites des approches existantes et
motivant les choix méthodologiques retenus dans ce travail.

Sur le plan expérimental, les trois modèles de base sont évalués selon un protocole
strict incluant une validation croisée, une gestion explicite du déséquilibre de classes
(environ 65\,\% de cas négatifs contre 35\,\% de cas positifs) avec Borderline-SMOTE,
ainsi qu'une transformation de puissance des variables. L'évaluation repose sur un 
ensemble complet de métriques cliniquement pertinentes : précision (\textit{accuracy}), 
rappel (\textit{recall}), F1-score et aire sous la courbe ROC (AUC-ROC).
Des méthodes ensemblistes (Soft Voting et Stacking) ont également été éprouvées.

Les résultats expérimentaux révèlent que \textbf{XGBoost} avec sa configuration par 
défaut obtient les meilleures performances globales, avec notamment un AUC-ROC de 
0.9531 et un F1-Score de 0.8571, confirmant la supériorité des méthodes de boosting 
par gradient. Les méthodes plus complexes telles que le Stacking n'apportent pas d'amélioration 
significative. Les \textbf{Forêts Aléatoires} présentent des performances très comparables,
tandis que les \textbf{SVM} avec noyau RBF offrent de bonnes performances mais peinent
à rivaliser sur le diagnostic précis des patients atteints. Mener l'optimisation des hyperparamètres 
ne s'est pas avéré essentiel vu la robustesse du feature engineering et des modèles.

Ces travaux contribuent à fournir une base comparative reproductible et cliniquement orientée
pour le choix d'un algorithme de détection du diabète, et ouvrent des perspectives vers des
approches de déploiement en système multi-agents au service de la santé.

\medskip
\noindent\textbf{Mots-clés :} Apprentissage automatique, Diabète, PIDD, SVM, Forêts
Aléatoires, XGBoost, Classification, AUC-ROC, Données médicales.